\sgaStableTarget{sga6.e13.sec.6}\subsection*{6. Appendix: Study of @@M0@@-sheaves on @@M1@@.}

\sgaStableTarget{sga6.e13.para.6-1}\textbf{6.1.} Let @@M2@@ be an algebraically closed field, and let @@M3@@. Let @@M4@@ be a coherent sheaf on @@M5@@, with Hilbert polynomial
@@M6@@
and let
@@M7@@
be a sequence of integers @@M8@@. Put
@@M9@@
Finally, for every integer @@M10@@, denote the largest coherent subsheaf of @@M11@@ with support of dimension @@M12@@ by @@M13@@, and put @@M14@@.

\sgaStableTarget{sga6.e13.lem.6-2}\textbf{Lemma 6.2.} \emph{Under the conditions of 6.1, consider an exact sequence}
@@M15@@
\emph{with @@M16@@. Then @@M17@@ is @@M18@@-regular with @@M19@@, where @@M20@@ is the @@M21@@-th @@M22@@-polynomial (1.9).}

Indeed, @@M23@@ is clearly a @@M24@@-sheaf (@@M25@@ terms), so @@M26@@ is a @@M27@@-sheaf by 1.6. By 2.10, @@M28@@. Thus the assertion follows from 1.11 and 1.10.

\sgaStableTarget{sga6.e13.prop.6-3}\textbf{Proposition 6.3.} \emph{Under the conditions of 6.1, suppose that @@M29@@ is a @@M30@@-sheaf. Then one can calculate, by means of universal polynomials in the @@M31@@, @@M32@@ and @@M33@@, integers @@M34@@ and @@M35@@ such that there is an exact sequence}
@@M36@@
\emph{Moreover, @@M37@@ occurs only in the calculation of @@M38@@.}

